%% kapa: exam=1 round=26 year=2015 period=1 form=A subject=civil_law subject_ko=민법 seq=14 number=14
%% source: Q-Net 감정평가사 기출문제 (cst00309) artlSeq=5138001 file=extracted/1008/5138001/제26회_감정평가사 1차_기출/제26회_감정평가사 1차_기출/1교시_감평_A형.pdf
%% conversion: kordoc 4.24.51 -> tools/gen1.py (auto); review=proofread-sonnet; flags=none
\begin{question}{14}
표현대리에 관한 설명으로 옳지 않은 것은? (다툼이 있으면 판례에 따름)
\begin{choices}
  \item 대리권 소멸 후의 표현대리가 인정되는 경우, 그 표현대리의 권한을 넘은 대리행위가 있을 때 권한을 넘은 표현대리가 성립할 수 있다.
  \item 대리인이 대리권 소멸 후 복대리인을 선임하여 복대리인으로 하여금 대리행위를 하도록 한 경우, 대리권 소멸 후의 표현대리가 성립할 수 없다.
  \item 권한을 넘은 표현대리에 있어 ‘정당한 이유’의 존부는 대리행위 당시의 제반 사정을 객관적으로 관찰하여 판단하여야 한다.
  \item 본인이 대리권 수여표시에 의한 표현대리 또는 대리권 소멸 후의 표현대리로 인하여 책임을 지기 위해서는 상대방이 선의ㆍ무과실이어야 한다.
  \item 교회의 대표자가 교인총회의 결의를 거치지 아니하고 교회 재산을 처분한 행위에 대하여 권한을 넘은 표현대리에 관한 규정을 준용할 수 없다.
\end{choices}
\end{question}
\answer{2}
